% Draft: engagement-definition section for the ToG extension (2026-09-08).
% Citation keys follow docs/references/definition_refs.bib plus the CoG paper's refer.tex
% (ke2022, tot2021camera, riotapi2024timeline). Numbers come from docs/DEFINITION_EVIDENCE.md;
% every figure is traceable to a script under features/fight_boundary/.

\section{Defining the Engagement}\label{sec:definition}

\subsection{The construct and its operationalisation}

``Teamfight'' is a player concept, not a measurement unit. Neither the publisher nor the
literature provides a formal definition: the official wiki describes a teamfight as
``multiple (e.g.\ the entire team) champions gather[ing] in one area to battle'', and the
formal definitions that exist -- Schubert et al.'s \emph{encounter}
\cite{schubert2016encounter}, Ke et al.'s \emph{team fight} \cite{ke2022}, and the
OpenDota labels used by Tot et al.\ \cite{tot2021camera} -- share one atom (champions of
both teams within range of each other at the same time) and differ in whether a kill is
required and how many champions must take part. We therefore treat the engagement as an
operationalised construct in the sense of measurement modelling
\cite{jacobs2021measurement}: a rule that maps observable telemetry to instances of an
unobservable concept, whose assumptions must be made explicit and \emph{tested} for
reliability (can it be repeated?) and validity (is it right?) rather than asserted. The
rest of this section states the rule, derives its constants from the corpus, and reports
the tests.

\subsection{Kill episodes and the temporal boundary}

Public Match-V5 telemetry \cite{riotapi2024timeline} records no damage events; the only
combat event with a millisecond timestamp and a position is the champion kill. The unit is
therefore built from kills, as in Ke et al.\ \cite{ke2022}, and the population is
explicitly the set of combat episodes that resolve with at least one death. This is a
\emph{case definition} in the epidemiological sense: encounters without a kill exist but
are outside the observable population. Schubert et al.\ report that even with complete
damage logs 81.1\% of proximity encounters contained no kill, half of those no damage at
all, and they restricted outcome analysis to encounters with kills
\cite{schubert2016encounter}; in our corpus, encounters with three or more alive champions
per side sustained for 20\,s but no kill amount to 4.9\% of engagements.

Consecutive kills of a match are chained into an episode when their gap is at most $G$.
Rather than fixing $G$ by hand we follow the practice established for session
identification from inter-activity times \cite{halfaker2015session,mehrzadi2012session}
and for earthquake declustering from nearest-neighbour space--time distances
\cite{zaliapin2013clusters}: on a logarithmic scale the distribution of inter-kill
intervals is bimodal, with a within-fight mode near 6\,s and a between-fight mode near
62\,s, and the boundary is placed at the antimode. Halfaker et al.\ locate the valley by
the crossing of a two-component Gaussian mixture; Mehrzadi and Feitelson search for the
deepest minimum inside a bounded range because global thresholds imprint artefacts on the
resulting session-length distribution. We combine the two: a mixture locates the
components, and the valley is the minimum of a kernel density estimate between them
\cite{silverman1981multimodality}. Across the full corpus of 208,141 matches
(10.4 million intervals) the valley is at 13.7\,s (match-level bootstrap 95\% CI
12.2--15.7, computed on 100,000-interval replicates and therefore conservative), stable
under a three-fold change of bandwidth (13.6--14.1\,s), across the three patches
(13.5--14.0\,s), and, in a 9,903-match subsample, across match intensity terciles
(11.4--14.5\,s). The mixture crossing gives 14.5\,s. Clustering agreement (adjusted Rand
index $\geq 0.9$) holds for $G \in [10, 18]$\,s, so the value is reported together with
this plateau, in the spirit of stability-based cluster selection
\cite{campello2013hdbscan}.

Two honest qualifications follow from the same analysis. First, the valley is a mid-game
feature: after 20 minutes the interval density is unimodal, a continuum from 3\,s to 60\,s
filled by chases, staggered deaths and respawn returns, so in the late game $G$ is a
convention justified only by the plateau. Second, Mehrzadi and Feitelson warn that using
activity gaps both to place boundaries and to judge them is circular
\cite{mehrzadi2012session}; the spatial boundary below is therefore validated against an
independent signal.

\subsection{The spatial boundary}

An episode is split when two of its kills are more than $D$ apart (complete linkage), so
that simultaneous fights on different parts of the map are not merged. $D$ is estimated
from an independent signal: whether the two kills involve the same champion (as killer,
victim or assister). Two kills within seconds that share a champion belong to one fight
by construction, and a champion cannot be in two places, so within the temporal window
the sharing rate as a function of distance traces where ``same fight'' ends. Its 50\%
crossing -- the Bayes decision point that misclassifies the fewest pairs -- is at
4,264 units (95\% CI 4,256--4,274) over the 5.46 million consecutive pairs within $G$;
it lies between 3,974 (top lane) and 4,412 units (bases) in every map region (lanes,
river, jungle, both pits, bases), within 7\% of the pooled value, so a single $D$
suffices. Pairs within 7\,s but farther than 4,000 units share a champion only 3--5\%
of the time: these are simultaneous fights elsewhere. The sharing rate is 93.6\% for pairs
the rule links and 10.9\% for the simultaneous-but-distant pairs it separates. Inside
lanes the boundary is anisotropic: holding the offset across the lane below 500 units,
the crossing along the lane lies at 5,200--5,900 units, whereas holding the along-lane
offset below 1,000 units the crossing across the lane lies at 3,450--3,730 units, a
ratio of about 1.5 that recurs in every patch. We keep the isotropic rule: an elliptical
boundary would relabel 0.16\% of in-window pairs, and the anisotropy survives when
walkable-path distance replaces Euclidean distance (below), so it describes how fights
move along a lane rather than the terrain.

The seismological scalar proximity $\eta = \Delta t\,\Delta d^{\,d}$
\cite{zaliapin2013clusters,zaliapin2022perspectives} applied to the same kill pairs is
bimodal for every $d$ we tried, and at $d=2$ its partition agrees with ours on 92\% of
pairs with a marginally lower separation of shared from unshared pairs (93.3\% vs 93.7\%
sharing among linked pairs at equal link rate); the rectangular rule is thus a slightly
sharper special case of an established space--time clustering criterion. Walkable-path
distance computed on the minimap's terrain mask did not improve the separation
(AUC 0.939 vs 0.940), and 91\% of consecutive pairs are not separated by walls; the
boundary is a property of how fights move, not of terrain.

\subsection{Presence gate, anchor and lead}

An episode becomes an \emph{engagement} when, at the cutoff $\tau$ placed $B$ before
its first kill, at least $M=2$ alive champions of each team stand within $R$ of the first
kill's position, mirroring Ke et al.'s requirement that a team fight involve at least two
players on a side \cite{ke2022}. Neither $R$ nor $B$ can be estimated from Match-V5,
whose positions are interpolated between 60\,s frames, so we take the two constants the
game itself uses for the same purpose. $R = 1{,}600$ units is the radius within which the
game credits an ally with proximity to a champion's death (it shares the death's
experience that far); $B = 15$\,s is the window within which the game credits a champion
with contribution to a kill, and it is also the offset at which OpenDota's team-fight
label, used by Tot et al.\ \cite{tot2021camera}, starts before the first death. Both are
rules rather than estimates: they apply to every champion and every patch, and a change
would appear in the patch notes. The bracket around them is narrow -- champions see each
other within 1,350 units, the longest non-ultimate engage ability (Zac's Elastic
Slingshot at rank five) reaches 1,800, and a champion leaves combat 5\,s after the last
exchange -- and $B$ is a lead, not a fight onset. Because the gate is a rule-anchored
convention, we report its leverage (Appendix~\ref{app:gate}). On 81,914 kill episodes
from 2,971 matches the share that qualifies as an engagement is 10.3\% at (1,600 units,
15\,s) and 19.1\% at the point used in our CoG paper (1,800 units, 10\,s), monotone in
both constants; the gate, not $G$ or $D$, decides how many engagements a match contains.
Re-running the tabular baseline at both points with every other constant fixed gives
531,984 engagements at a held-out test AUC of 0.660 (95\% CI 0.656--0.663) against
994,365 at 0.669 (0.665--0.672), and on the 32,820 test engagements both gates accept
the two models are indistinguishable (0.658 vs 0.659). The gate changes how many
engagements there are, not how predictable they are.
Moving the anchor from the first kill to the episode's kill centroid (median shift 663
units) changes the verdict for 6\% of episodes with no directional bias in presence
counts or scale classes.

\subsection{Re-estimation per patch}

Because esport rules change with every patch and models built on one iteration of the
game decay on the next \cite{chitayat2023beyond}, the constants are not fixed once. A
pipeline re-estimates $G$, $D$ and the coverage of $R$ for each patch and applies a
pre-registered drift rule: a patch keeps the pooled definition if its $G$ lies inside the
pooled plateau, its $D$ within 10\% of the pooled value with overlapping intervals, and the
game rules behind $R$ and $B$ unchanged in its patch notes. For patches 15.14--15.16 the per-patch values are
$G = 14.0, 13.5, 13.7$\,s and $D = 4{,}285, 4{,}265, 4{,}241$ units (within 0.5\% of the
pooled value), all inside, and one definition serves the corpus. A pilot on 3,000 matches
per patch had placed patch 15.15 at 12.6\,s and 4,146 units; on the full corpus that
difference vanishes, so the tolerances of the drift rule are wide relative to the
between-patch variation actually observed, and a genuine change of pace would be
unmistakable. This mirrors the stability analysis of Zaliapin and
Ben-Zion, who report their cluster measure moving within $\pm0.1$ when catalogue
parameters are varied by orders of magnitude \cite{zaliapin2013clusters}.

\subsection{Scale classes}

Scale is a post-hoc attribute (participation is known only when the fight resolves) and is
used only for reporting. Its classes follow modes of the data, as behavioural states are
read from density peaks in unsupervised ethology \cite{berman2014mapping}: on the diagonal
of the per-team participation distribution 2v2 (12.2\%) and 5v5 (8.1\%) are modes with a
trough at 4v4, and the block with four or more on each side is 22.4\% of engagements. The
smaller side's count alone has no antimode (21.9, 36.0, 19.7, 14.2, 8.1\% for 1--5), so
the cut between two and three per side is a convention and is labelled as such. We use
\emph{pick} (smaller side $\leq 1$; 83\% of these are one champion caught by two or three),
\emph{skirmish} (2--3) and \emph{teamfight} ($\geq 4$, matching the community sense of
``the entire team'').

\subsection{Validity summary}

Following \cite{jacobs2021measurement}: content validity -- the rule instantiates the
shared atom of all published definitions and states what it excludes; convergent validity
-- it coincides with Ke et al.'s team fight and with the seismological proximity
partition on 92\% of pairs; discriminant validity -- linked pairs share a champion at 94\%
against 51\% for separated pairs; test--retest reliability -- three patches, bootstrap
intervals and the ARI plateau; predictive validity -- with every input restricted to what is
observable at the cutoff, the engagement outcome is predictable from pre-cutoff state at
AUC 0.67 (0.69 under the minute-gold currency), and this predictability does not depend
on the scale of the engagement; the scale gradient we reported earlier was carried by a
time feature normalised by the match's own length (Section~\ref{sec:results}). Construct validity is a
matter of degree; the pipeline and its verdicts are released so that the degree can be
re-assessed on new patches.
